% Part: first-order-logic
% Chapter: natural-deduction
% Section: quantifier-rules

\documentclass[../../../include/open-logic-section]{subfiles}

\begin{document}

\olfileid{fol}{ntd}{qrl}

\olsection{د سورونو قاعدې}

\subsection{د $\lforall$ دپاره قاعدې}

\begin{defish}
\AxiomC{$!A(a)$}
\RightLabel{\Intro{\lforall}}
\UnaryInfC{$\lforall[x][\Atom{!A}{x}]$}
\DisplayProof
\hfill
\AxiomC{$\lforall[x][\Atom{!A}{x}]$}
\RightLabel{\Elim{\lforall}}
\UnaryInfC{$!A(t)$}
\DisplayProof
\end{defish}

د~$\lforall$ په قاعدو کښې $t$ يوه تړلې اصطلاح ده (يعنې داسې اصطلاح
چې هېڅ متغير نۀ لري)، او $a$~داسې !!a{constant} دے چې په
پايله~$\lforall[x][!A(x)]$ يا په هرې هغې فرضيې کښې نۀ راځي چې د
مقدمې~$!A(a)$ په پای ته رسېدونکي !!{derivation} کښې
!!{undischarged} پاتې وي۔ موږ $a$ د~\Intro{\lforall} استنتاج
\emph{ځانګړے متغير} بولو۔\footnote{موږ د “ځانګړي متغير” اصطلاح
کاروو، که څه هم په پورتنۍ قاعده کښې $a$ يو ثابت دے۔ دا نوم تاريخي
لاملونه لري۔}

\subsection{د $\lexists$ دپاره قاعدې}

\begin{defish}
\AxiomC{$\Atom{!A}{t}$}
\RightLabel{\Intro{\lexists}}
\UnaryInfC{$\lexists[x][\Atom{!A}{x}]$}
\DisplayProof
\hfill
\AxiomC{$\lexists[x][\Atom{!A}{x}]$}
\AxiomC{[$\Atom{!A}{a}$]$^n$}
\DeduceC{$!C$}
\DischargeRule{\Elim{\lexists}}{n}
\BinaryInfC{$!C$}
\DisplayProof
\end{defish}

بيا هم، $t$ يوه تړلې اصطلاح ده، او $a$ داسې !!a{constant} دے چې
په مقدمه~$\lexists[x][!A(x)]$، پايله~$!C$، يا په هرې هغې فرضيې
کښې نۀ راځي چې د دواړو مقدمو په پای ته رسېدونکو !!{derivation}s کښې
!!{undischarged} پاتې وي (د~$!A(a)$ فرضيو پرته)۔ موږ $a$ د
\Elim{\lexists} استنتاج \emph{ځانګړے متغير} بولو۔

هغه شرط چې په~\Intro{\lforall} يا~\Elim{\lexists} استنتاج کښې
ځانګړے متغير په پايله يا په هرې هغې فرضيې کښې نۀ راځي چې مقدمو ته
په رسېدونکو !!{derivation}s کښې !!{undischarged} وي، او د
په وجودي مقدمه کښې هم نۀ راځي، \emph{د ځانګړي
متغير شرط} بلل کېږي۔ \textbf{د ژباړې د نسخې سمون (OLND-002):}
سرچينه په دې عمومي جمله کښې وايي چې ځانګړے متغير په “مقدمو” کښې
نۀ راځي، خو د دواړو قاعدو مرستندويه مقدمه همدا ځانګړے ثابت لري؛
دلته د سرچينې له مخکښې کره شرطونو سره سم پايله، نۀ ايستل شوې فرضيې
او د وجود ايستلو وجودي مقدمه ټاکل شوې دي۔

\begin{explain}
هغه دود راياد کړئ چې کله $!A$ داسې !!a{formula} وي چې
!!{variable}~$x$ پکې ازاد وي، نو دا په~$!A(x)$ ليکلو ښيو۔ په همدې
چوکاټ کښې بيا $!A(t)$ د~$\Subst{!A}{t}{x}$ لنډون دے۔ نو د
$\Intro\lexists$ قاعده داسې هم ليکلے شو:
\begin{prooftree}
\AxiomC{$\Subst{!A}{t}{x}$}
\RightLabel{\Intro{\lexists}}
\UnaryInfC{$\lexists[x][!A]$}
\end{prooftree}
پام وکړئ چې $t$ کېدے شي له مخکښې په~$!A$ کښې راغلے وي؛ مثلاً
$!A$~کېدے شي~$\Atom{\Obj P}{t,x}$ وي۔ له دې امله، له
~$\Atom{\Obj P}{t,t}$ څخه $\lexists[x][\Atom{\Obj P}{t,x}]$
استنتاجول د~$\Intro\lexists$ سمه کارونه ده---تاسو د مقدمې د~$t$
يوه يا زياتې پېښې په تړلي !!{variable}~$x$ “بدلولے” شئ، او اړينه
نۀ ده چې ټولې پېښې ئې بدلې کړئ۔ خو د~$\Intro\lforall$ او
~$\Elim\lexists$ د ځانګړي متغير شرطونه غواړي چې !!{constant}~$a$
په~$!A$ کښې نۀ راځي۔ نو د~$\Intro\lforall$ په کارولو سره له
$\Atom{\Obj P}{a,a}$ څخه $\lforall[x][\Atom{\Obj P}{a,x}]$ په سمه
توګه نۀ شئ استنتاجولے۔
\end{explain}

\begin{explain}
په~\Intro{\lexists} او~\Elim{\lforall} کښې له تړلتيا پرته بل
بنديز نشته، او اصطلاح~$t$ هره تړلې اصطلاح کېدے شي؛ نو د بل شرط
اندېښنه نۀ شته۔ بل لوري ته، د~\Elim{\lexists} او
\Intro{\lforall} په قاعدو کښې د ځانګړي متغير شرط غواړي چې
!!{constant}~$a$ په پايله يا کومې !!{undischarged} فرضيې کښې هېڅ
ځاے نۀ راځي۔ دا شرط د نظام د سموالي د يقيني کولو دپاره اړين دے؛
يعنې نظام له هغو !!{undischarged} فرضيو څخه يوازې هغه !!{sentence}s
!!{derive}s کړي چې ترې معنايي پايله کېږي۔ له دې شرط پرته لاندې به
روا و:
\begin{prooftree}
  \AxiomC{$\lexists[x][!A(x)]$}
  \AxiomC{$\Discharge{!A(a)}{1}$}
  \RightLabel{*\Intro{\lforall}}
  \UnaryInfC{$\lforall[x][!A(x)]$}
  \RightLabel{\Elim{\lexists}}
  \BinaryInfC{$\lforall[x][!A(x)]$}
\end{prooftree}
خو $\lexists[x][!A(x)] \Entails/ \lforall[x][!A(x)]$۔

ځکه چې د سورونو د ايستلو قاعدې د !!{variable}s پر ځاے يوازې د تړلو
اصطلاحاتو ايښوولو اجازه ورکوي، نو هره هغه !!{formula} چې د
!!{sentence}s له يوه سټ څخه وايستل شي، پخپله !!a{sentence} ده۔
\end{explain}


\end{document}
